\chapter{The Investigation Mode}
\label{ch:investigation}

\begin{versebox}
What question was asked? What answer was given?\\
And what verse sums it up?\\
Investigate the sutta --- that's this mode.
\end{versebox}

\noindent If the Teaching Mode shows you the \textit{structure} of a lesson, the Investigation Mode shows you how to \textit{take it apart}. What exactly does this mode investigate? Everything.\footnote{Pali: \textit{``Tattha katamo vicayo hāro? `Yaṃ pucchitañca vissajjitañcā'ti gāthā, ayaṃ vicayo hāro.''} The Netti answers its own question by pointing back to the mnemonic verse from Chapter~3. Then comes the list of what gets investigated.}

The terms. The questions. The answers. What comes before and after. The appeal, the danger, the escape, the fruit, the means, the command, the summary verse --- and all nine suttas that serve as the Netti's test cases.\footnote{Pali: \textit{``Kiṃ vicinati? Padaṃ vicinati, pañhaṃ vicinati, visajjanaṃ vicinati, pubbāparaṃ vicinati, assādaṃ vicinati, ādīnavaṃ vicinati, nissaraṇaṃ vicinati, phalaṃ vicinati, upāyaṃ vicinati, āṇattiṃ vicinati, anugītiṃ vicinati, sabbe nava suttante vicinati.''} Twelve objects of investigation. Note that the first six components of the Teaching Mode (appeal through command) are all included --- the Investigation Mode doesn't replace the Teaching Mode, it \textit{incorporates} it. It also adds three more: the term (\textit{pada}), the before-and-after (\textit{pubbāpara}), and the follow-up verse (\textit{anugīti}).}

To demonstrate, the Netti picks one of the most haunting exchanges in the entire Pali Canon: the Venerable Ajita's questions to the Buddha from the Pārāyana --- the ``Way to the Far Shore'' --- the final and most ancient section of the Suttanipāta.\footnote{The Pārāyaṇa (Sn 5) is widely regarded as one of the oldest strata of the Pali Canon. It consists of sixteen sets of questions posed by brahmin students to the Buddha, each in verse form with verse answers. The Netti uses the first set --- Ajita's questions (Sn 5.1, \textit{Ajitamāṇavapucchā}) --- as its demonstration text for the Investigation Mode. The choice is significant: these verses are cryptic, compressed, and multi-layered, making them ideal for showing how the Investigation Mode cracks open a difficult text.}

\section*{Ajita's First Question: What's Wrong with the World?}

\begin{versebox}
``By what is the world wrapped up?\\
Why does it not shine forth?\\
What do you call its smearing?\\
What is its greatest fear?''
\end{versebox}

Four questions in four lines. But the Netti says: actually, this is \textit{one} question with four aspects. Why? Because all four lines are about the same thing --- the world.\footnote{Pali: \textit{``Kenassu nivuto loko, kenassu nappakāsati; / Kissābhilepanaṃ brūsi, kiṃ su tassa mahabbhaya''nti.''} (Sn 5.1, v. 1038). The Netti's analysis: \textit{``Imāni cattāri padāni pucchitāni, so eko pañho. Kasmā? Ekavatthu pariggahā''} --- ``These four terms were asked; it is one question. Why? Because they all grasp a single topic.'' Ajita asks about the world's covering, its failure to shine, its smearing, and its great fear --- four aspects of one diagnosis.}

The word ``world'' here has three meanings: the world of defilements, the world of existence, and the world of the senses.\footnote{Pali: \textit{``Loko tividho kilesaloko bhavaloko indriyaloko.''} This threefold analysis of ``world'' (\textit{loka}) is standard in the Pali commentarial tradition but the Netti uses it specifically to show that Ajita's question --- and the Buddha's answer --- operates simultaneously on all three levels. When the Buddha says ``the world is covered by ignorance,'' he means all three worlds: the psychological world of defilements, the cosmological world of rebirth, and the experiential world of the senses.}

Now the Buddha answers --- and watch how each line of the answer matches its question, one to one:

\begin{versebox}
``By ignorance the world is wrapped, Ajita,''\\
said the Blessed One.\\
``Through meanness and heedlessness it doesn't shine.\\
Craving I call its smearing.\\
Suffering is its greatest fear.''
\end{versebox}

The Netti unpacks each pair:\footnote{Pali: \textit{``Avijjāya nivuto loko, ajitāti bhagavā, vivicchā pamādā nappakāsati; / Jappābhilepanaṃ brūmi, dukkhamassa mahabbhaya''nti.''} (Sn 5.1, v. 1039). The Netti then says: \textit{``Imāni cattāri padāni imehi catūhi padehi visajjitāni paṭhamaṃ paṭhamena, dutiyaṃ dutiyena, tatiyaṃ tatiyena, catutthaṃ catutthena''} --- ``These four terms were answered by these four terms: first by first, second by second, third by third, fourth by fourth.'' The Investigation Mode's first job is to establish this precise correspondence between question and answer.}

\textbf{``By what is the world wrapped?'' --- By ignorance.} Not just any ignorance. The Netti identifies it as the master hindrance: ``All beings are obstructed by the hindrance of ignorance.'' The Buddha himself says: ``I declare that there is, in a sense, just one hindrance for all beings --- namely, ignorance.''\footnote{Pali: \textit{``Nīvaraṇehi nivuto loko, avijjānīvaraṇā hi sabbe sattā. Yathāha bhagavā `sabbasattānaṃ, bhikkhave, sabbapāṇānaṃ sabbabhūtānaṃ pariyāyato ekameva nīvaraṇaṃ vadāmi yadidaṃ avijjā, avijjānīvaraṇā hi sabbe sattā.'\,''} The cross-reference is powerful: the Netti pulls in the Buddha's own declaration that ignorance is the single root hindrance. This is the Synonym Mode in action within the Investigation Mode --- recognizing that ``the world is covered'' and ``all beings are obstructed by ignorance'' are two ways of saying the same thing.}

\textbf{``Why doesn't it shine?'' --- Through doubt and heedlessness.} When you're covered by ignorance, you doubt. When you doubt, you don't commit. When you don't commit, you don't put in effort. When you don't put in effort, wholesome qualities don't arise. And what doesn't arise can't shine.\footnote{Pali: \textit{``Yo puggalo nīvaraṇehi nivuto, so vivicchati. Vivicchā nāma vuccati vicikicchā.''} Here the Netti gives an etymology: \textit{vivicchā} (meanness, reluctance) is identified with \textit{vicikicchā} (doubt, uncertainty). The causal chain: ignorance → doubt → no faith → no effort → no wholesome qualities → no shining. The Netti then quotes Dhp 304: ``The good shine from afar, like the Himalaya mountains. The bad are not seen here, like arrows shot in the night.''}

\textbf{``What is its smearing?'' --- Craving.} The Pali word \textit{jappā} (yearning, craving) is identified as another name for \textit{taṇhā}. Craving ``smears'' the world the way grease smears glass --- it coats everything and distorts your vision.\footnote{Pali: \textit{``Jappā nāma vuccati taṇhā. Sā kathaṃ abhilimpati?''} --- ``\textit{Jappā} is a name for craving. How does it smear?'' The Netti then quotes: ``The person stained by passion doesn't know what's beneficial, doesn't see the Dhamma. Utter darkness then prevails when passion overcomes a person.'' (A citation also found at It 84). The image of craving as a \textit{coating} that prevents clear seeing is one of the Netti's most vivid metaphors.}

\textbf{``What is its greatest fear?'' --- Suffering.} All beings recoil from suffering. There is no fear comparable to it, let alone greater.\footnote{Pali: \textit{``Sabbe sattā hi dukkhassa ubbijjanti, natthi bhayaṃ dukkhena samasamaṃ, kuto vā pana tassa uttaritaraṃ?''} --- ``All beings are terrified of suffering; there is no fear equal to suffering, let alone surpassing it.'' The Netti then distinguishes three types of suffering: \textit{dukkhadukkhatā} (the suffering of pain itself), \textit{vipariṇāmadukkhatā} (the suffering of change), and \textit{saṅkhāradukkhatā} (the suffering inherent in all conditioned things). The first two, says the Netti, are occasionally escaped --- some beings live long and healthy lives. But the third is escaped only by final nibbāna. Therefore it is \textit{saṅkhāradukkhatā} that is the world's ``greatest fear.''}

\section*{Ajita's Second Question: What Stops the Flood?}

\begin{versebox}
``Streams flow everywhere ---\\
what restrains those streams?\\
Tell me what blocks them:\\
by what are those streams shut off?''
\end{versebox}

Now Ajita shifts from diagnosis to cure. If ignorance, doubt, craving, and suffering are the problem --- what's the solution?\footnote{Pali: \textit{``Savanti sabbadhi sotā, sotānaṃ kiṃ nivāraṇaṃ; / Sotānaṃ saṃvaraṃ brūhi, kena sotā pidhīyare.''} (Sn 5.1, v. 1040). The Netti analyzes these as two questions, not one, because they ask for two different things: ``what restrains?'' asks for a temporary check, while ``what blocks?'' asks for a permanent closure. This distinction between suppression (\textit{vikkhambhana}) and uprooting (\textit{samugghāta}) is fundamental to the Abhidhamma.}

The Netti says these are actually \textit{two} questions. ``What restrains the streams?'' asks for the temporary check. ``What blocks them?'' asks for the permanent shutdown.

\begin{versebox}
``Whatever streams there are in the world, Ajita,''\\
said the Blessed One,\\
``mindfulness is their restraint.\\
I tell you the blocking of the streams:\\
by wisdom they are shut off.''
\end{versebox}

\textbf{Mindfulness restrains.} When mindfulness focused on the body is developed and cultivated, the eye doesn't get swept away by pleasant sights or crash against unpleasant ones. The ear doesn't. The nose doesn't. And so on through all six senses.\footnote{Pali: \textit{``Kāyagatāya satiyā bhāvitāya bahulīkatāya cakkhu nāviñchati manāpikesu rūpesu, amanāpikesu na paṭihaññati, sotaṃ\ldots ghānaṃ\ldots jivhā\ldots kāyo\ldots mano nāviñchati.''} The verb \textit{āviñchati} means ``to pull, to drag'' --- the untrained senses pull the mind toward pleasant objects and slam against unpleasant ones. Mindfulness prevents both. Note the emphasis on \textit{kāyagatā sati} (mindfulness directed to the body) --- the Netti specifies body-based mindfulness, not mindfulness in general, as the restraining factor.}

\textbf{Wisdom uproots.} Mindfulness dams the river. But wisdom pulls out the underlying tendencies by the root --- ``like uprooting the trunk of a great tree so that the succession of flowers, fruits, leaves, and shoots is cut off entirely.''\footnote{Pali: \textit{``Paññāya anusayā pahīyanti, anusayesu pahīnesu pariyuṭṭhānā pahīyanti.''} --- ``By wisdom the underlying tendencies are abandoned; when the underlying tendencies are abandoned, the obsessions are abandoned.'' Then the simile: \textit{``Taṃ yathā khandhavantassa rukkhassa anavasesamūluddharaṇe kate pupphaphalapallavaṅkurasantati samucchinnā bhavati.''} --- ``Just as when a great tree's root is completely pulled up, the succession of flowers, fruits, leaves, and shoots is cut off.'' \textit{Anusaya} (underlying tendency, latent disposition) lies deeper than \textit{pariyuṭṭhāna} (obsession, active defilement). Mindfulness handles the surface waves; wisdom deals with the undertow.}

This is a critical distinction. Mindfulness handles the surface. Wisdom handles the depth. You need both, but they do different things.

\section*{Ajita's Third Question: Where Does Everything Cease?}

\begin{versebox}
``Wisdom and mindfulness,\\
name-and-form, dear sir ---\\
tell me, when asked: where do these cease?''
\end{versebox}

Ajita has heard that mindfulness restrains and wisdom uproots. Now he asks the ultimate question: where does the whole process come to an end? Where do wisdom, mindfulness, and even name-and-form --- the entirety of psychophysical existence --- finally stop?\footnote{Pali: \textit{``Paññā ceva sati ca, nāmarūpañca mārisa; / Etaṃ me puṭṭho pabrūhi, katthetaṃ uparujjhatī''ti.''} (Sn 5.1, v. 1042). \textit{Nāmarūpa} (name-and-form) in early Buddhism means the totality of experienced reality: \textit{nāma} (the mental factors --- feeling, perception, intention, contact, attention) and \textit{rūpa} (the physical forms --- the five sense faculties). Ajita is asking: when does \textit{everything} stop?}

The Buddha's answer is one of the most important verses in the Canon:

\begin{versebox}
``This question you have asked, Ajita ---\\
I will answer you.\\
Where name-and-form ceases\\
without remainder ---\\
with the cessation of consciousness,\\
right here it ceases.''
\end{versebox}

With the cessation of consciousness, name-and-form ceases.\footnote{Pali: \textit{``Yametaṃ pañhaṃ apucchi, ajita taṃ vadāmi te; / Yattha nāmañca rūpañca, asesaṃ uparujjhati; / Viññāṇassa nirodhena, etthetaṃ uparujjhatī''ti.''} (Sn 5.1, v. 1043). This directly parallels the dependent origination formula: \textit{viññāṇapaccayā nāmarūpaṃ} --- ``consciousness conditions name-and-form.'' When consciousness ceases, name-and-form ceases. The Netti elaborates: \textit{``Nāmarūpañca viññāṇahetukaṃ viññāṇapaccayā nibbattaṃ, tassa maggena hetu upacchinno, viññāṇaṃ anāhāraṃ anabhinanditaṃ appaṭisandhikaṃ taṃ nirujjhati''} --- ``Name-and-form has consciousness as its cause, arises from consciousness as its condition. By the path, the cause is cut off. Consciousness, unfed, undelighted-in, without re-linking --- that ceases.''}

The Netti unpacks this with characteristic precision. Name-and-form is caused by consciousness. The path cuts that cause. When consciousness is no longer ``fed'' by craving, no longer ``delighted in,'' no longer capable of ``re-linking'' to a new existence --- it ceases. And with it, all of name-and-form ceases.\footnote{The three descriptions of consciousness's cessation are technically precise: \textit{anāhāra} (unfed --- cutting off the ``nutriment'' of consciousness, i.e., craving), \textit{anabhinandita} (undelighted-in --- the end of the delight that sustains rebirth), and \textit{appaṭisandhi} (without relinking --- no more connecting to a new existence). These map to the three ``rounds'' (\textit{vaṭṭa}) of dependent origination: the round of defilements, the round of action, and the round of results.}

\section*{Ajita's Fourth Question: The Conduct of the Liberated}

\begin{versebox}
``Those who have fully understood,\\
and the many trainees here ---\\
wise one, when asked, tell me their conduct.''
\end{versebox}

Three groups: the fully realized (arahants), the trainees (those on the path but not yet finished), and --- implicitly --- those just beginning insight practice.\footnote{Pali: \textit{``Ye ca saṅkhātadhammāse, ye ca sekkhā puthū idha; / Tesaṃ me nipako iriyaṃ, puṭṭho pabrūhi mārisā''ti.''} (Sn 5.1, v. 1044). The Netti analyzes this as three questions: \textit{saṅkhātadhamma} (those who have fully understood = arahants), \textit{sekkhā} (trainees = stream-enterers through once-returners), and the implied question about \textit{vipassanāpubbaṅgamappahānayoga} (the practice of abandonment preceded by insight = the pre-path practice).}

The Buddha's answer:

\begin{versebox}
``One should not be greedy for sense pleasures,''\\
said the Blessed One to Ajita.\\
``One should be unmuddied in mind.\\
Skilled in all phenomena ---\\
a monk should wander, mindful.''
\end{versebox}

Four instructions, and the Netti maps each one to a specific practice:\footnote{Pali: \textit{``Kāmesu nābhigijjheyya, manasānāvilo siyā; / Kusalo sabbadhammānaṃ, sato bhikkhu paribbaje''ti.''} (Sn 5.1, v. 1045). The Netti's detailed analysis: ``not greedy for sense pleasures'' = guarding the senses against craving for pleasant objects; ``unmuddied in mind'' = overcoming obsessive defilements through the four right efforts; ``skilled in all phenomena'' = full comprehension (\textit{pariññā}) through seeing and development; ``mindful'' = clear comprehension in all activities.}

\begin{itemize}[itemsep=4pt]
\item \textbf{``Not greedy''} --- Guard the senses. When craving for pleasant objects arises, don't follow it. The trainee has to actively protect the mind against two things: greed for attractive objects and aversion toward unattractive ones.\footnote{Pali: \textit{``Sekho dvīsu dhammesu cittaṃ rakkhitabbaṃ gedhā ca rajanīyesu dhammesu, dosā ca pariyuṭṭhānīyesu.''} --- ``The trainee must guard the mind regarding two things: greed toward attractive phenomena, and aversion toward irritating phenomena.''}

\item \textbf{``Unmuddied in mind''} --- The four right efforts. Prevent unwholesome states from arising. Abandon those that have arisen. Generate wholesome states. Maintain those already present.\footnote{The Netti maps each of the three unwholesome thought-types to a specific faculty that counteracts it: \textit{kāmavitakka} (sensual thought) is countered by the \textit{satindriya} (mindfulness faculty), \textit{byāpādavitakka} (ill-will thought) by the \textit{samādhindriya} (concentration faculty), and \textit{vihiṃsāvitakka} (cruel thought) by the \textit{vīriyindriya} (energy faculty). The \textit{paññindriya} (wisdom faculty) then eliminates all unwholesome states that have already arisen. Faith (\textit{saddhindriya}) is the confidence in these four faculties.}

\item \textbf{``Skilled in all phenomena''} --- Full comprehension. Knowing what's wholesome and unwholesome, blameworthy and blameless, dark and bright, what to cultivate and what to avoid.\footnote{Pali: \textit{``Kusalo sabbadhammānaṃ\ldots''} The Netti unpacks ``skilled'' (\textit{kusala}) into eight kinds of skill: skilled in meaning (\textit{atthakusala}), skilled in phenomena (\textit{dhammakusala}), skilled in what's beneficial (\textit{kalyāṇatākusala}), skilled in results (\textit{phalatākusala}), skilled in what leads forward (\textit{āyakusala}), skilled in what leads astray (\textit{apāyakusala}), skilled in method (\textit{upāyakusala}), and endowed with great skill (\textit{mahatā kosallena samannāgata}). This is comprehensive competence.}

\item \textbf{``Mindful''} --- Clear comprehension in every activity. Going forward, going back, looking ahead, looking aside, bending, stretching, wearing robes, carrying the bowl, eating, drinking, walking, standing, sitting, sleeping, waking, speaking, and being silent.\footnote{This is the \textit{sampajañña} (clear comprehension) formula from DN 22 and MN 10, the Satipaṭṭhāna Suttas. The Netti quotes it here to show that ``mindfulness'' in the Buddha's final answer to Ajita isn't vague attentiveness --- it's a specific, comprehensive practice of awareness applied to every bodily activity throughout the day.}
\end{itemize}

\section*{How to Investigate Any Sutta}

The chapter ends with a remarkable paragraph that gives you the meta-instructions: how to apply the Investigation Mode to \textit{any} sutta you encounter.\footnote{Pali: \textit{``Evaṃ pucchitabbaṃ, evaṃ visajjitabbaṃ. Suttassa ca anugīti atthato ca byañjanato ca samānetabbā.''} --- ``Thus should questions be asked, thus should they be answered. And the sutta's follow-up verse should be connected both in meaning and in phrasing.'' Then comes the key instruction: meaning without proper phrasing is ``idle chatter'' (\textit{samphappalāpa}). Badly placed words make even the meaning hard to follow. Therefore the teaching must be recited with \textit{both} meaning and phrasing intact.}

When you sit down with a sutta, ask these questions:

\begin{itemize}[itemsep=4pt]
\item Is this sutta a direct statement, or does it build on something said earlier?
\item Is its meaning explicit, or does it need to be drawn out?
\item Does it deal with the contaminated side (defilements and their causes), or the penetrating side (insight and its results)?
\item Is it at the level of a trainee, or of one beyond training?
\item Where in this sutta can all four truths be seen --- at the beginning, the middle, or the end?
\end{itemize}

That's the Investigation Mode: take the sutta apart into its question, its answer, and its summary verse. Then investigate every element. Not just what was said, but why it was said, what question prompted it, and how the answer maps onto the whole system.\footnote{The closing line: \textit{``Evaṃ suttaṃ pavicetabbaṃ. Tenāha āyasmā mahākaccāyano --- `yaṃ pucchitañca vissajjitañca, suttassa yā ca anugītī'ti.''} --- ``Thus should a sutta be investigated. Therefore said the Venerable Mahākaccāyana: `What was asked and what was answered, and what is the sutta's follow-up verse.'\,'' The section concludes as it began, returning to the mnemonic verse, now fully unpacked.}

\ornament

\noindent Next: the Fitness Mode --- the quality check. How do you know if your reading of a sutta is actually valid?
